% GENERATED FILE. Edit canonical lessons, then run npm run generate:books.
% canonical-source-hash: fnv1a64:5c52d1cf93229f4a
% canonical-lessons: JA-C275-tsunami, JA-C275-nagareru, JA-C275-nagasu, JA-C275-uku, JA-C275-mata

\chapter{The Great Wave}
\label{ch:ja-the-great-wave}
\hlchaptermodality{275}

\begin{chapteropening}
\textbf{By the end of this chapter:} \emph{I can say what the sea carried away.}

Complete the last lesson of chapter 275: I can say what the sea carried away.
\end{chapteropening}

\section[\texorpdfstring{\ja{つなみ} (tsunami)}{つなみ (tsunami)}]{\ja{つなみ} (tsunami) — a tsunami; a tidal wave}

\label{lesson:JA-C275-tsunami}



\begin{quote}
Before the new one: say the Japanese for frightening, then the Japanese for evacuation.
\end{quote}

\subsection*{You'll want to know: \ja{つなみ}}
\textbf{\ja{つなみ}} — \emph{tsunami} — \textquotedblleft{}a tsunami; a tidal wave\textquotedblright{}.

\begin{quote}
\textbf{\ja{つなみの} \ja{ばあいは、やまに} \ja{にげて} \ja{ください。}} — \emph{tsunami no baai wa, yama ni nigete kudasai} — \textquotedblleft{}In case of a tsunami, please run to high ground.\textquotedblright{}
\end{quote}

\begin{cousinweb}
\textbf{\ja{つなみ}} (\emph{tsunami}) is a native Japanese word (wago), inherited rather than borrowed from Chinese, built from \emph{tsu}, \textquotedblleft{}harbour\textquotedblright{}, and \emph{nami}, \textquotedblleft{}wave\textquotedblright{} (JMdict, KANJIDIC2).
\end{cousinweb}

\begin{grammarlens}[title={What you've built}]
One more word for telling what happened.
\end{grammarlens}

\subsection*{Your turn}
\begin{itemize}
\raggedright
  \item \emph{Say it:} \emph{tsunami}
  \item \emph{Say it:} \emph{tsunami}, once more
  \item \emph{Say it:} \emph{tsunami no baai wa, yama ni nigete kudasai} — In case of a tsunami, please run to high ground.
  \item \emph{Recall:} say the Japanese for frightening, then the Japanese for evacuation, then say \emph{tsunami} again
\end{itemize}

\begin{tcolorbox}[breakable,colback=teal!4,colframe=teal!35!black,title={Before you move on}]
What does \ja{つなみ} mean? (\textquotedblleft{}A tsunami; a tidal wave\textquotedblright{}.) Say it once more.
\end{tcolorbox}
\section[\texorpdfstring{\ja{ながれる} (nagareru)}{ながれる (nagareru)}]{\ja{ながれる} (nagareru) — to flow; to be washed away}

\label{lesson:JA-C275-nagareru}



\begin{quote}
Before the new one: say the Japanese for evacuation, then the Japanese for a tsunami.
\end{quote}

\subsection*{You'll want to know: \ja{ながれる}}
\textbf{\ja{ながれる}} — \emph{nagareru} — \textquotedblleft{}to flow; to be washed away\textquotedblright{}.

\begin{quote}
\textbf{\ja{かわの} \ja{みずが} \ja{はやく} \ja{ながれています。}} — \emph{kawa no mizu ga hayaku nagarete imasu} — \textquotedblleft{}The river is running fast.\textquotedblright{}
\end{quote}

Polite \textbf{\ja{ながれます}}, \emph{nagaremasu}; an ichidan verb. Past (polite): \textbf{\ja{くるまが} \ja{ながれました。}}, \emph{kuruma ga nagaremashita}.

\begin{cousinweb}
\textbf{\ja{ながれる}} (\emph{nagareru}) is a native Japanese word (wago), inherited rather than borrowed from Chinese; its kanji was attached later as a kun'yomi (JMdict).
\end{cousinweb}

\begin{grammarlens}[title={What you've built}]
One more word for telling what happened.
\end{grammarlens}

\subsection*{Your turn}
\begin{itemize}
\raggedright
  \item \emph{Say it:} \emph{nagareru}
  \item \emph{Say it:} \emph{nagareru}, once more
  \item \emph{Say it:} \emph{kawa no mizu ga hayaku nagarete imasu} — The river is running fast.
  \item \emph{Recall:} say the Japanese for evacuation, then the Japanese for a tsunami, then say \emph{nagareru} again
\end{itemize}

\begin{tcolorbox}[breakable,colback=teal!4,colframe=teal!35!black,title={Before you move on}]
What does \ja{ながれる} mean? (\textquotedblleft{}To flow; to be washed away\textquotedblright{}.) Say it once more.
\end{tcolorbox}
\section[\texorpdfstring{\ja{ながす} (nagasu)}{ながす (nagasu)}]{\ja{ながす} (nagasu) — to wash away; to let flow}

\label{lesson:JA-C275-nagasu}



\begin{quote}
Before the new one: say the Japanese for a tsunami, then the Japanese for to flow.
\end{quote}

\subsection*{You'll want to know: \ja{ながす}}
\textbf{\ja{ながす}} — \emph{nagasu} — \textquotedblleft{}to wash away; to let flow\textquotedblright{}.

\begin{quote}
\textbf{\ja{つなみが} \ja{うちを} \ja{ながしました。}} — \emph{tsunami ga uchi o nagashimashita} — \textquotedblleft{}The tsunami swept our house away.\textquotedblright{}
\end{quote}

Polite \textbf{\ja{ながします}}, \emph{nagashimasu}; a godan verb. Past (polite): \textbf{\ja{つなみが} \ja{くるまを} \ja{ながしました。}}, \emph{tsunami ga kuruma o nagashimashita}.

\begin{cousinweb}
\textbf{\ja{ながす}} (\emph{nagasu}) is a native Japanese word (wago), inherited rather than borrowed from Chinese; its kanji was attached later as a kun'yomi (JMdict). It is the transitive partner of \emph{nagareru}, \textquotedblleft{}to flow\textquotedblright{}.
\end{cousinweb}

\begin{grammarlens}[title={What you've built}]
One more word for telling what happened.
\end{grammarlens}

\subsection*{Your turn}
\begin{itemize}
\raggedright
  \item \emph{Say it:} \emph{nagasu}
  \item \emph{Say it:} \emph{nagasu}, once more
  \item \emph{Say it:} \emph{tsunami ga uchi o nagashimashita} — The tsunami swept our house away.
  \item \emph{Recall:} say the Japanese for a tsunami, then the Japanese for to flow, then say \emph{nagasu} again
\end{itemize}

\begin{tcolorbox}[breakable,colback=teal!4,colframe=teal!35!black,title={Before you move on}]
What does \ja{ながす} mean? (\textquotedblleft{}To wash away; to let flow\textquotedblright{}.) Say it once more.
\end{tcolorbox}
\section[\texorpdfstring{\ja{うく} (uku)}{うく (uku)}]{\ja{うく} (uku) — to float}

\label{lesson:JA-C275-uku}



\begin{quote}
Before the new one: say the Japanese for to flow, then the Japanese for to wash away.
\end{quote}

\subsection*{You'll want to know: \ja{うく}}
\textbf{\ja{うく}} — \emph{uku} — \textquotedblleft{}to float\textquotedblright{}.

\begin{quote}
\textbf{\ja{みずの} \ja{うえに} \ja{きが} \ja{ういていました。}} — \emph{mizu no ue ni ki ga uite imashita} — \textquotedblleft{}Wood was floating on the water.\textquotedblright{}
\end{quote}

Polite \textbf{\ja{うきます}}, \emph{ukimasu}; a godan verb. Past (polite): \textbf{\ja{きが} \ja{うきました。}}, \emph{ki ga ukimashita}.

\begin{cousinweb}
\textbf{\ja{うく}} (\emph{uku}) is a native Japanese word (wago), inherited rather than borrowed from Chinese; its kanji was attached later as a kun'yomi (JMdict).
\end{cousinweb}

\begin{grammarlens}[title={What you've built}]
One more word for telling what happened.
\end{grammarlens}

\subsection*{Your turn}
\begin{itemize}
\raggedright
  \item \emph{Say it:} \emph{uku}
  \item \emph{Say it:} \emph{uku}, once more
  \item \emph{Say it:} \emph{mizu no ue ni ki ga uite imashita} — Wood was floating on the water.
  \item \emph{Recall:} say the Japanese for to flow, then the Japanese for to wash away, then say \emph{uku} again
\end{itemize}

\begin{tcolorbox}[breakable,colback=teal!4,colframe=teal!35!black,title={Before you move on}]
What does \ja{うく} mean? (\textquotedblleft{}To float\textquotedblright{}.) Say it once more.
\end{tcolorbox}
\section[\texorpdfstring{\ja{また} (mata)}{また (mata)}]{\ja{また} (mata) — once more; again}

\label{lesson:JA-C275-mata}



\begin{quote}
Before the new one: say the Japanese for to wash away, then the Japanese for to float.
\end{quote}

\subsection*{You'll want to know: \ja{また}}
\textbf{\ja{また}} — \emph{mata} — \textquotedblleft{}once more; again\textquotedblright{}.

\begin{quote}
\textbf{\ja{また} \ja{つなみが} \ja{くるでしょう。}} — \emph{mata tsunami ga kuru deshō} — \textquotedblleft{}Another tsunami will probably come.\textquotedblright{}
\end{quote}

\begin{cousinweb}
\textbf{\ja{また}} (\emph{mata}) is a native Japanese word (wago), inherited rather than borrowed from Chinese, and it is usually written in kana (JMdict).
\end{cousinweb}

\begin{grammarlens}[title={What you've built}]
One more word for telling what happened.
\end{grammarlens}

\subsection*{Your turn}
\begin{itemize}
\raggedright
  \item \emph{Say it:} \emph{mata}
  \item \emph{Say it:} \emph{mata}, once more
  \item \emph{Say it:} \emph{mata tsunami ga kuru deshō} — Another tsunami will probably come.
  \item \emph{Recall:} say the Japanese for to wash away, then the Japanese for to float, then say \emph{mata} again
\end{itemize}

\begin{tcolorbox}[breakable,colback=teal!4,colframe=teal!35!black,title={Before you move on}]
What does \ja{また} mean? (\textquotedblleft{}Once more; again\textquotedblright{}.) Say it once more.
\end{tcolorbox}
